\chapter{УДИРТГАЛ}

\section{Судалгааны үндэслэл, ач холбогдол}

3D реконструкц нь робот техник, виртуал (VR) болон AR, соёлын өвийн дижитал архив, автономи систем зэрэг олон салбарт үндсэн технологийн үүрэгтэй. 2D зурагнуудаас 3D загвар үүсгэх нь хүний оролцоогүйгээр 3D загварчлал хийх, 3D хувьчлахад чухал ач холбогдолтой.

Уламжлалт 3D сэргээн босголтын аргууд болох Structure-from-Motion (SfM) болон Multi-View Stereo (MVS) нь бодит геометрийг өндөр нарийвчлалтайгаар сэргээдэг боловч дараах сул талуудтай:

\begin{itemize}
    \item \textbf{Гэрэлтүүлгийн хувирлаас өндөр хамааралтай}
    \item \textbf{Объектын гадаргуугийн тусгал, хээ, өнгөнөөс шалтгаалан алдаа үүсдэг}
    \item \textbf{Олон зураг шаарддаг}
    \item \textbf{Тооцооллын хувьд хүнд}
\end{itemize}

Сүүлийн жилүүдэд хиймэл оюун ухааны хөгжлийн ачаар 2D зургаас 3D бүтцийг сэргээх гүн сургалтын (Deep Learning) аргууд үүссэн. Гэсэн хэдий ч эдгээр аргууд нь их хэмжээний өгөгдөл шаарддаг, бодит цагийн хэрэглээнд удаан, заримдаа бүтэц алдагддаг зэрэг сул талуудтай.

Тиймээс энэхүү судалгааны ажлын гол зорилго нь:

\begin{enumerate}
    \item Уламжлалт SfM + MVS аргууд
    \item Орчин үеийн гүн сургалтын аргууд (CNN/Point Cloud methods)
\end{enumerate}

хоёрын боломж, гүйцэтгэлийг ижил орчинд харьцуулж, **хэмжилтийн тогтвортой нөхцөл** бүрдүүлэх камерын төхөөрөмж (controlled capture rig) ашиглан 3D сэргээн босголтын чанарыг сайжруулах арга замыг тодорхойлох явдал юм.

\section{Судалгааны зорилго, зорилт}

\subsection{Зорилго}

2D зурагнуудаас 3D загвар үүсгэх уламжлалт болон AI суурьтай аргуудын гүйцэтгэлийг харьцуулан судалж, тогтвортой гэрэлтүүлэгтэй зураг авах орчин бүтээж, реконструкцийн чанарыг сайжруулахад оршино.

\subsection{Зорилтууд}

\begin{enumerate}
    \item \textbf{Зураг авах орчин зохиох:} Гэрэлтүүлэг тогтмол, нэгэн ижил нөхцөлд зураг авах боломжтой камерын жижиг орчин (mini capture box) зохион бүтээх
    \item \textbf{Өгөгдөл цуглуулах:} Объектуудыг олон өнцгөөс, жигд гэрэлтүүлэгтэйгээр зураг авах
    \item \textbf{Уламжлалт аргуудыг хэрэгжүүлэх:} COLMAP ашиглан SfM ба MVS реконструкц хийх
    \item \textbf{Гүн сургалтын аргуудыг турших:} CNN ба point-based reconstruction model-уудыг турших
    \item \textbf{Гүйцэтгэлийг харьцуулах:}
    \begin{itemize}
        \item Chamfer Distance
        \item Point-to-surface error
        \item Visual geometric accuracy
    \end{itemize}
    \item \textbf{Автоматжуулсан pipeline бүтээх:} Зураг оруулахад SfM → MVS → Mesh generation алхмуудыг нэг даралтаар ажиллуулдаг скрипт боловсруулах
\end{enumerate}

\section{Судалгааны хамрах хүрээ}

Судалгааны ажлын хүрээ дараах байдлаар хязгаарлагдана:

\begin{itemize}
    \item \textbf{Объект:} Жижиг хэмжээтэй, нэг ширхэг хатуу биет объект
    \item \textbf{Зураг авах орчин:} Тогтмол гэрэлтүүлэг бүхий хаалттай камерын төхөөрөмж
    \item \textbf{Арга:} SfM + MVS болон AI reconstruction аргуудын харьцуулалт
    \item \textbf{Програм хангамж:} COLMAP, OpenMVS, Python, PyTorch
    \item \textbf{Гаралт:} Point cloud болон mesh хэлбэртэй 3D загвар
\end{itemize}

\section{Судалгааны шинэлэг тал}

Энэхүү судалгааны ажил дараах өвөрмөц онцлог, шинэлэг талтай:

\begin{enumerate}
    \item \textbf{Тогтмол гэрэлтүүлэг бүхий камерын орчин} — реконструкцийн тогтвортой байдлыг нэмэгдүүлэх зорилготой өөрийн бүтээсэн capture rig ашигласан
    \item \textbf{Уламжлалт ба AI аргуудын шууд харьцуулалт} — ижил өгөгдөлд хоёр аргын чанарын ялгааг нарийвчлан үнэлсэн
    \item \textbf{Автомат pipeline} — зураг авалтаас mesh гаралт хүртэл бүх алхмыг автоматжуулсан
    \item \textbf{Mesh чанарын тоон ба дүрслэх үнэлгээ} — зөвхөн харааны үнэлгээ бус Chamfer distance ашигласан
    \item \textbf{Бодит хэрэглээнд шууд ашиглаж болох шийдэл} — жижиг лабораторид, музей, соёлын өвийн сканнерт ашиглаж болох прототип
\end{enumerate}
